\par\Needspace{6\baselineskip}\medskip
\hypertarget{native-114-article085}{}\pdfbookmark[2]{85.}{article-085}
\begin{center}85.\end{center}
Retinebimus characteres $r$, $r'$, $N$, $N'$, $p$, $e$, $\Pi$ in eadem significatione, in qua supra accepti sunt; differentiam $N'-N$ denotabimus per $\Delta$, tempusque intra quod corpus coeleste a loco priori ad posteriorem movetur per $t$. Iam patet, si valor approximatus alicuius quantitatum $p$, $e$, $\Pi$ sit notus, etiam duas reliquas inde determinari posse, ac dein per methodos in Sectione prima explicatas tempus motui a loco primo ad secundum respondens. Quod si tempori proposito $t$ aequale evadit, valor suppositus ipsius $p$, $e$ vel $\Pi$ est ipse verus, orbitaque ipsa iam inventa; sin minus, calculus cum valore alio a primo parum diverso repetitus docebit, quanta variatio in valore temporis variationi exiguae in valore ipsius $p$, $e$, $\Pi$ respondeat, unde per simplicem interpolationem valor correctus eruetur. Cum quo si calculus denuo repetitur, tempus emergens vel ex asse cum proposito quadrabit, vel saltem perparum ab eo differet, ita ut certe novis correctionibus adhibitis consensum tam exactum attingere liceat, quantum tabulae logarithmicae et trigonometricae permittunt.

Problema itaque eo reductum est, ut pro eo casu, ubi orbita adhuc penitus incognita est, valorem saltem approximatum alicuius quantitatum $p$, $e$, $\Pi$ \SourceLeaf{115}determinare doceamus. Methodum iam trademus, per quam valor ipsius $p$ tanta praecisione eruitur, ut pro parvis quidem valoribus ipsius $\Delta$ nulla amplius correctione indigeat, adeoque tota orbita per primum calculum omni iam praecisione determinetur, quam tabulae vulgares permittunt. Vix unquam autem aliter nisi pro valoribus mediocribus ipsius $\Delta$ ad hanc methodum recurrere oportebit, quum determinationem orbitae omnino adhuc incognitae, propter problematis complicationem nimis intricatam, vix aliter suscipere liceat, nisi per observationes non nimis ab invicem distantes, aut potius tales, quibus motus heliocentricus non nimius respondet.
